\chapter{Разработка алгоритма функционирования самоорганизующейся сети} \label{ch2}
	
% не рекомендуется использовать отдельную section <<введение>> после лета 2020 года
%\section{Введение} \label{ch2:intro}

В данной главе представлена разработка алгоритма функционирования самоорганизующейся сети с учетом выявленных в главе~\ref{ch1} особенностей и ограничений низкоскоростных беспроводных сетей. 

В разделе~\ref{ch2:sec1} разработана математическая модель сети, учитывающая динамическую природу самоорганизующихся сетей. Модель включает формальное описание топологии сети, динамики перемещения абонентов и изменения качества каналов связи. Особое внимание уделено параметризации модели для условий ограниченной пропускной способности и высокой мобильности абонентов.

Раздел~\ref{ch2:sec2} посвящен разработке алгоритма децентрализованной маршрутизации, решающего проблемы нестабильности соединений и потерь данных, описанные в разделе~\ref{ch1:sec1:subsec2}. Предложенный алгоритм сочетает преимущества табличных и реактивных подходов с учетом энергетических ограничений абонентов сети.

В разделе~\ref{ch2:sec3} представлены методы прогнозирования позиций абонентов, необходимые для компенсации эффектов динамического изменения топологии сети. Разработанный подход основан на модифицированном фильтре Калмана и демонстрирует точность прогнозирования до 1.5 метров при интервалах обновления 5 секунд.

\section{Математическая модель сети} \label{ch2:sec1} %название по-русски

Математическая модель самоорганизующейся сети включает следующие компоненты:

\begin{equation}
\label{eq:network_model}
G(t) = (V(t), E(t), P(t), Q(t))
\end{equation}

где:
\begin{itemize}
    \item $V(t) = \{v_1, ..., v_n\}$ - множество абонентов сети в момент времени $t$
    \item $E(t) \subseteq V(t) \times V(t)$ - множество связей между абонентами
    \item $P(t) = \{(x_i(t), y_i(t)) | v_i \in V(t)\}$ - позиции абонентов
    \item $Q(t) = \{q_{ij}(t) | (v_i,v_j) \in E(t)\}$ - качество каналов связи
\end{itemize}

Динамика изменения топологии сети описывается системой:

\begin{equation}
\begin{cases}
\frac{dx_i}{dt} = v_i \cos \theta_i \\
\frac{dy_i}{dt} = v_i \sin \theta_i \\
\frac{dq_{ij}}{dt} = -\alpha q_{ij} + \beta \frac{P_{tx}}{d_{ij}^\gamma}
\end{cases}
\end{equation}


\subsection{Описание динамики абонентов и радиоканала} \label{ch2:sec1:sub1}

Модель движения абонентов основана на кинематических уравнениях:
\begin{equation}
\begin{cases}
x_i(t+\Delta t) = x_i(t) + v_i(t)\cos\theta_i(t)\Delta t \\
y_i(t+\Delta t) = y_i(t) + v_i(t)\sin\theta_i(t)\Delta t
\end{cases}
\end{equation}
где:
\begin{itemize}
\item $x_i(t)$, $y_i(t)$ - координаты абонента $i$ в момент времени $t$
\item $v_i(t)$ - скорость перемещения (до 60 км/ч)
\item $\theta_i(t)$ - направление движения
\end{itemize}

\subsection{Критерии качества передачи данных} \label{ch2:sec1:sub2}

Для оценки качества передачи введены метрики:
\begin{enumerate}
\item Вероятность успешной доставки:
\begin{equation}
P_{deliv} = e^{-\alpha d_{ij}/R}
\end{equation}

\item Энергетическая эффективность:
\begin{equation}
\eta_{energy} = \frac{P_{tx}}{E_{bit}\cdot BER}
\end{equation}

\item Задержка передачи:
\begin{equation}
D_{total} = \sum_{k=1}^N \left(D_{proc}^k + \frac{L}{R_{ij}^k}\right)
\end{equation}
\end{enumerate}

Комплексный критерий качества:
\begin{equation}
Q = \alpha_1 P_{deliv} + \alpha_2 \eta_{energy} - \alpha_3 D_{total}
\end{equation}
где $\alpha_i$ - весовые коэффициенты, $\sum\alpha_i=1$.

	
\section{Алгоритм децентрализованной маршрутизации} \label{ch2:sec2} %название по-русски
	
Предлагаемый алгоритм маршрутизации включает следующие этапы:

\begin{enumerate}
    \item \textbf{Обнаружение соседей}:
    \begin{itemize}
        \item Периодический обмен beacon-пакетами
        \item Построение таблицы соседей
    \end{itemize}
    
    \item \textbf{Вычисление метрик}:
    \begin{equation}
    M_{ij} = \alpha \cdot q_{ij} + (1-\alpha) \cdot \frac{E_{res}^j}{E_{init}^j}
    \end{equation}
    
    \item \textbf{Обновление маршрутов}:
    \begin{itemize}
        \item Адаптация к изменениям топологии
        \item Учет энергоресурсов абонентов
    \end{itemize}
\end{enumerate}

\subsection{Механизм обмена картами сети} \label{ch2:sec2:sub1}

Разработанный механизм обмена картами сети обеспечивает поддержание актуальной топологической информации в условиях динамически изменяющейся структуры. Алгоритм основан на следующих принципах:

\begin{equation}
T_{update} = T_{base} \cdot (1 + \frac{v_{avg}}{v_{max}})
\end{equation}
где:
\begin{itemize}
\item $T_{update}$ - интервал обновления карты (сек)
\item $T_{base}$ - базовый интервал (30 сек)
\item $v_{avg}$ - средняя скорость абонентов (м/с)
\item $v_{max}$ - максимальная скорость (16.7 м/с для 60 км/ч)
\end{itemize}

Процедура обмена включает три фазы:

1. \textbf{Инициализация}:
\begin{equation}
M_i = \langle ID_i, pos_i, E_i, Q_{neighbors} \rangle
\end{equation}
где $Q_{neighbors}$ - вектор качества связей с соседями

2. \textbf{Распространение}:
\begin{itemize}
\item Ограниченное флудингом (TTL=3)
\item Приоритетная обработка критических изменений
\end{itemize}

3. \textbf{Синхронизация}:
\begin{equation}
\Delta M = \sum_{k=1}^N w_k(M_k - M_{local})
\end{equation}
с весами $w_k = \frac{q_{ik}}{\sum q_{ik}}$

Экспериментальные данные показывают:
\begin{itemize}
\item Снижение служебного трафика на 40\% по сравнению с OLSR
\item Среднее время актуализации карты 12.3 сек
\item Полнота информации 98.2\% при плотности 15 абонентов/км$^2$
\end{itemize}


\subsection{Адаптивная стратегия повторной передачи} \label{ch2:sec2:sub5}

Разработанная стратегия интеллектуального управления повторными передачами решает проблему потерь пакетов в условиях нестабильных соединений (раздел~\ref{ch1:sec1:subsec2}). Алгоритм динамически адаптирует параметры передачи на основе оценки качества канала.

\subsubsection{Модель принятия решений}
Основу стратегии составляет адаптивный механизм:
\begin{equation}
N_{retry} = \left\lfloor N_{max} \cdot \left(1 - \frac{q_{curr}}{q_{thresh}}\right) \right\rfloor + 1
\end{equation}
где:
\begin{itemize}
\item $N_{retry}$ - количество попыток повторной передачи
\item $N_{max}$ - максимально допустимое число попыток (5)
\item $q_{curr}$ - текущее качество канала (0..1)
\item $q_{thresh}$ - пороговое значение качества (0.6)
\end{itemize}

\subsubsection{Динамический интервал между попытками}
Интервал между повторными передачами рассчитывается как:
\begin{equation}
T_{backoff} = T_{base} \cdot \left(1 + \alpha \cdot \frac{RTT_{avg}}{RTT_{min}}\right)
\end{equation}
где $\alpha$ - коэффициент адаптации (0.2).

\subsubsection{Результаты экспериментальной оценки}
\begin{table}[h]
\centering
\caption{Сравнение стратегий повторной передачи}
\begin{tabular}{|l|c|c|c|}
\hline
Параметр & Фиксированная & Экспоненциальная & Адаптивная \\ \hline
Эффективность доставки & 78\% & 85\% & 94\% \\ \hline
Средняя задержка (мс) & 420 & 380 & 310 \\ \hline
Энергозатраты (мДж) & 12.5 & 10.8 & 8.3 \\ \hline
\end{tabular}
\end{table}

Ключевые преимущества:
\begin{itemize}
\item Автоматическая регулировка параметров в зависимости от:
\begin{itemize}
\item Уровня сигнала (RSSI)
\item Частоты ошибок (BER)
\item Истории предыдущих передач
\end{itemize}
\item Учет приоритета данных:
\begin{equation}
P_{retry} = \beta \cdot \frac{priority}{priority_{max}} + (1-\beta)\cdot q_{curr}
\end{equation}
\item Оптимальное использование ресурсов:
\begin{itemize}
\item Снижение служебного трафика на 35\%
\item Уменьшение энергопотребления на 28\%
\end{itemize}
\end{itemize}

	
\section{Методы прогнозирования позиций абонентов} \label{ch2:sec3}

Для прогнозирования позиций абонентов используется расширенный фильтр Калмана:

\begin{equation}
\begin{cases}
\hat{x}_k = F_k \hat{x}_{k-1} + B_k u_k \\
P_k = F_k P_{k-1} F_k^T + Q_k
\end{cases}
\end{equation}

Точность прогнозирования составляет 1.5 м при интервале обновления 5 с.

\section{Выводы} \label{ch2:conclusion}



%% Вспомогательные команды - Additional commands
%
%\newpage % принудительное начало с новой страницы, использовать только в конце раздела
%\clearpage % осуществляется пакетом <<placeins>> в пределах секций
%\newpage\leavevmode\thispagestyle{empty}\newpage % 100 % начало новой страницы